%%%%

%%%   Самарский государственный технический университет

%%%   Кафедра прикладной математики и информатики

%%%

%%%   Преамбула для статьи в журнал "Вестник Самарского государственного

%%%   технического университета. Серия: «Физико-математические науки»"

%%%   Ориентирована под MikTEx 2.4 for Win32

%%%

%%%   Хотя сам журнал собирается под GNU/Linux на дистрибутиве TeXLive



\documentclass[11pt,a4paper,twoside]{article}



\usepackage[cp1251]{inputenc}

\usepackage[T2A]{fontenc}

\usepackage[english,russian]{babel}

\usepackage{samgtubib}

\usepackage[dvips]{graphicx}

\usepackage[multiple]{footmisc}



\usepackage{samgtu}

\renewcommand{\VestnikYear}{2009} % Год издания Вестника

\renewcommand{\VestnikNumber}{2\,(19)} % Номер Вестника

\renewcommand{\VestnikMonth}{Ноябрь} % Месяц выхода Вестника

\renewcommand{\VestnikData}{01.11.09} % Дата подписания Вестника в печать

\renewcommand{\VestnikUPL}{26,64} % Усл. печ. л.

\renewcommand{\VestnikUIL}{25,73} % Уч.-изд. л.

\renewcommand{\VestnikRegNumber}{479/09} % Регистрационный номер

\renewcommand{\VestnikOrder}{994} % Номер заказа







  \def\pad{{\mathfrak p}}

%% %
%  \begin{document}
  
\renewcommand{\baselinestretch}{.95}

\renewcommand{\firstpage}{251}
\renewcommand{\lastpage}{259}


 % Номер УДК
 \udk{517.9}%
 \msc{34G10; 28C20, 11D88}% Коды MSC см. http://www.ams.org/msc/

 \title{Функциональный  оператор Лапласа на $\pad$-адическом пространстве и~формулы  Фейнмана и~Фейнмана--Каца}% Полный заголовок
 {Функциональный оператор Лапласа на $\pad$-адическом пространстве \dots}% Заголовок, который идёт в верхний колонтитул

 \author{Н.\,Н.~Шамаров} {Ш\,а\,м\,а\,р\,о\,в~Н.\,Н.}

 \institute{Московский государственный университет им.~М.~В.~Ломоносова,\\
119992, Москва, Воробьёвы горы.}

 \email{E-mail}{nshamarov@yandex.ru} %E-mail's авторов

 \otherinf{\fullauthor{Николай Николаевич Шамаров}\regalia{к.ф.-м.н.}\position{научный сотрудник} \dept{институт теоретических проблем микромира им.~Н.~Н.~Боголюбова}}

 \keywords{формулы Фейнмана--Каца, формулы Фейнмана, функциональный лапласиан, $p$"=адический анализ}

 \workabstract{Строятся однородные замкнутые псевдодифференциальные операторы, аналогичные степеням бесконечномерного лапласиана,
 действующие в~банаховых пространствах комплексных функций,  определённых, в~свою очередь, на функциональных пространствах над полем $\mathfrak p$"=адических чисел. Для элементов полугрупп,  порождённых этими псевдодифференциальными операторами как генераторами, находятся формулы Фейнмана и~Фейнмана"--~Каца.}

 \engtitle{Functional Laplace operator on a  $\pad$"=adic space and Feynman--Kac and Feynman formulas} 

 \engauthor{N.\,N.~Shamarov}{S\,h\,a\,m\,a\,r\,o\,v~N.\,N.}


 \enginstitute{ M.~V.~Lomonosov Moscow State University, \\
Vorobjovy Gory, Moscow,  119899, Russia.}% Место работы каждого из авторов на англ. языке
 % Если несколько, то так  \enginstitute{ Организация 1 \and Организация 2  \and Организация 3}

 \otherenginfm{\fullauthor{Nikolaj~N. Shamarov}
 \regalia{Ph.\,D. (Phys. \& Math.)}
 \position{Scientific Fellow}
 \dept{N.~N.~Bogoliubov Institute for Theoretical Problems of Microphysics}
 }

\engabstract{Homogeneous closed PDO are constructed which are analogous to the powers of (absolute value of) infinite dimensional Laplacian and acting in Banach spaces of complex-valued functions defined on function spaces over a field of $\mathfrak p$-adic numbers. For elements of semigroups, for which these PDOs are generators, Feynman formulas and Feynman--Kac ones are obtained.}

 \engkeywords{Feynman--Kac formulas,
 Feynman formulas,
 functional Laplacian,
 $p$-adic analysis}

 \date{25/XII/2010} % Дата поступления статьи в редакцию.
 \revisiondate{17/V/2011} % В окончательном виде

 %%%%%%%%%%%%%%%%%%%%%%%%%%%%%%%%%%%%%%%%%%%%%%%%%%%%%%%%%%%%%%%%%%%%%%%%%%%%%%%%%%%%


 \maketitle

\Section[N]{Введение} Некоторые бесконечномерные аналоги оператора Лапласа  естественно возникают в~квантовой теории поля, отвечая кинетическому члену в~гамильтониане.
Другого рода бесконечномерные аналоги оператора Лапласа  используются в~уравнениях, эквивалентных уравнениям Янга"--~Миллса.
При этом областями определения функций, на которые действуют операторы, являются вещественные многообразия.

С другой стороны, трудности, возникающие при описании с~помощью  известных теорий, оперирующих  вещественными координатами,
процессов  на планковских масштабах,  привели к~идее рассмотреть $\pad$"=адический анализ  в~качестве альтернативного математического аппарата~[\citen{shamarov:bib:VVZ}].

В аппарате теории комплексных функций конечного числа $\pad$"=адических аргументов, то есть пробегающих поле $\pad$"=адических чисел Гензеля,
важную роль играет аналог оператора Лапласа,  определяемый как псевдодифференциальный оператор (ПДО)  и~называемый оператором Владимирова. 
Поскольку известен иной оператор, также носящий имя  В.~С.~Владимирова, то описанный выше ПДО Владимирова,  аналогичный оператору Лапласа, называем далее
 оператором Лапласа"--~Владимирова.


Изложенное выше мотивирует изучение аналогов оператора Лапласа"--~Владимирова,  действующих в пространстве комплекснозначных функций  бесконечномерного $\pad$"=адического аргумента.
Область определения этих функций, являющаяся бесконечномерным  $\pad$"=адическим пространством, называется далее конфигурационным  пространством. 
Пространство $\mathbb Q_\pad^{\mathbb N}$ счётных последовательностей $\pad$"=адических чисел в~качестве  конфигурационного рассмотрено как модельное в работах~[\citen{shamarov:bib:SSK}, \citen{shamarov:bib:BeloshQN}].
Данная же работа посвящена случаю функционального конфигурационного пространства,  которое представляется в~контексте упомянутых выше физических идей более естественным.

Формулой Фейнмана~[\citen{shamarov:bib:Stekl}]\ % в конфигурационном пространстве
называется представление решения задачи Коши эволюционного уравнения  в~виде предела последовательности кратных интегралов, в~которой эта кратность неограниченно растёт.
Формулой Фейнмана в~конфигурационном пространстве называется формула Фейнмана, в~которой кратные интегралы берутся по декартовым степеням области определения начальной функции  задачи (которая и~является конфигурационным пространством).
Если конфигурационное пространство является векторным (соотв., многообразием), то соответствующим фазовым пространством называется  произведение конфигурационного пространства на двойственное к~ нему (соотв., кокасательное расслоение конфигурационного многообразия),
 и~формулой Фейнмана в фазовом пространстве  называется формула Фейнмана, в~которой кратные интегралы берутся по декартовым степеням этого фазового пространства.
 Формулами Фейнмана"--~Каца в~конфигурационном (или в~фазовом) пространстве для той же задачи называется представление  решения задачи в виде интеграла по пространству траекторий (= отображений отрезка  со значениями)  в~соответствующем пространстве (по счётно"=аддитивной мере или по псевдомере, например, псевдомере Фейнмана).
 Ниже, если явно не оговорено обратное, под формулами Фейнмана и~Фейнмана"--~Каца будут пониматься эти формулы в~конфигурационном пространстве.

 В настоящей работе строятся однородные замкнутые ПДО,	 аналогичные степеням бесконечномерного лапласиана  (ср.~[\citen{shamarov:bib:S82}, \citen{shamarov:bib:AS}]  %\cite{S82},\cite{AS}
  для вещественного случая),  действующие в банаховых пространствах комплексных функций,  определённых, в~свою очередь, на функциональных пространствах 
 над полем \linebreak $\pad$"=адических чисел,  и находятся формулы Фейнмана и Фейнмана"--~Каца  (в~конфигурационном пространстве)  для элементов полугрупп,
 порождённых этими ПДО как генераторами.


Классический однородный эллиптический дифференциальный оператор второго порядка с~постоянными коэффициентами  (например, оператор Лапласа) от
 комплексной функции нескольких  вещественных переменных можно определить как
 псевдодифференциальный оператор, % (ПДО),
 символ которого равен квадратичной форме  от <<импульсных>> переменных, сопряжённых исходным.
 Степени (абсолютной величины) оператора Лапласа определяются  аналогично с~помощью символов, равных  соответствующим удвоенным
 степеням нормы вектора,  составленного из <<импульсных>> координат.
В~случае комплексных функций конечного числа переменных,  пробегающих поле не вещественных, но $\pad$-адических чисел, та же схема со степенями нормы
 приводит к определению операторов  Владимирова $D^\alpha$, являющихся аналогами степени $\alpha/2$ абсолютной величины оператора Лапласа.
В обоих случаях итоговый оператор % схема % приводит
получается в~результате % эквивалентно
последовательного применения трёх операций: обратного преобразования Фурье, умножения на  однородную функцию"=символ  от <<импульсных>> переменных
 и, наконец, прямого преобразования Фурье.
Данный алгоритм можно применить~и в~случае функций бесконечного числа переменных с~той  лишь разницей, что (обобщённые) меры,  получаемые после первого (обратного)  преобразования Фурье, на которые умножается символ,  невозможно отождествить с~функциями,  даже обобщёнными,  в~силу отсутствия (по теореме Вейля)  достаточно хорошей  меры Хаара на бесконечномерных  пространствах.

Бесконечномерные ПДО в~пространствах функций,  определённых на вещественных локально выпуклых пространствах,  были определены с~помощью описанного  способа О.~Г.~Смоляновым~[\citen{shamarov:bib:S82}].

\smallskip

\Section{Предварительные сведения и обозначения}
Поле $\pad$"=адических$^1$%
\footnotetext{$^1$Чтобы далее использовать
  стандартную пару переменных
  $(q, p)$ при определении ПДО,
  для простых чисел (в термине
  $p$-адическое число, например) выбрано иное,
  чем стандартный курсив, написание.
  Оно выбрано готическим, поскольку
  так иногда в литературе обозначаются % из тех соображений,  что
  максимальные идеалы целых чисел,
  определяемые простыми числами.
  } 
\!(\mbox{$\pad\in\{2, 3,$} $5, \dots\}$ --- простое натуральное число% (2, 3, 5, \dots )
  )
  чисел, обозначаемое $\mathbb Q_\pad$, %~\cite{VVZ}
  представляет собой  пополнение   поля %~$\Qnu$
  рациональных чисел относительно единственного нормирования$^2$%
\footnotetext{$^2$Нормированием на алгебраическом поле $(K, {+}_K, 0_K, \, {\cdot}_K, 1_K)$ называется всякая неотрицательная функция $N:K\to{\mathbb R}$
  такая, что $\forall k, l \in K$ ($N(k)=0 \iff k=0_K$),  \mbox{$N(k\,{\cdot}_K\,l)=$} $=N(k)N(l)$ и $N(k\,{+}_K\,l)\leqslant  N (k)+N(l)$\,).
  При этом локально равномерно непрерывными относительно   метрики $\rho_N(k, l)=N(l-k)$ и~по совокупности аргументов  являются соответственно отображения $N$ и ${+}_K:K\times K\to K$ и ${\cdot}_K:K\times K\to K$,   продолжения которых на $\rho_N$-пополнение $K_N$ поля $K$ превращают
  это пополнение снова в нормированное поле.
  }\!\!, %%%%%%%%%%%%%%\footnote
  переводящего $\pad$ в $\pad^{-1}$ и~всякое простое натуральное $\mathfrak q{\ne}\pad$ в~$1$.
 Продолжение этого нормирования на $\mathbb Q_\pad$ называется $\pad$"=адическим, и~его значения на  всяком $x\in\mathbb Q_\pad$ обозначается $|x|_\pad$.
    Векторные пространства над $\mathbb Q_\pad$ далее
  также называются  $\pad$"=адическими.

  Без ссылок формулируются далее факты из~[\citen{shamarov:bib:VVZ}], относящиеся к анализу ком\-плекс\-но\-знач\-ных функций  $\pad$"=адических аргументов.


Каждое $\pad $-адическое число $a$ допускает однозначное представление в~виде  (сходящегося в $\mathbb Q_\pad$)
   ряда 
$$ a=\mathop\sum\limits_{k\in\mathbb Z,\ k \geqslant -\log _\pad |a|_\pad} % ^{+\infty} 
  a_k\cdot \pad^{+k}=[a]_\pad+\{a\}_\pad,
$$ 
  где  $a_k\in\{0, 1, \dots, \pad-1\}$  для всякого  $  k\in\mathbb Z  $, причём для $a\ne 0$,    $\log _\pad |a|_\pad\in\mathbb Z$  
  и~$ a_{\log _\pad |a|_\pad}\ne0$  (${\log _\pad 0= -\infty},$ $a_{-\infty}=0$),
  $[a]_\pad=\sum\limits_{  k\geqslant 0} a_k\cdot \pad^{+k}$,
  и сумма
  $\{a\}_\pad= \sum\limits_{ -\log_\pad |a|_\pad \leqslant k<0 }  a_k\cdot \pad^{+k}$  конечна.
   При этом $\{a\}_\pad$  всегда является рациональным числом  из вещественного полуинтервала $[0; 1)$.



Непрерывный гомоморфизм ${\chi_{{}_{\mathbb Q_\pad}}} $ аддитивной группы поля $\mathbb Q_\pad$  в~мультипли\-кативную группу $ \{z\in\mathbb C:|z|=1\} $,
задаваемый формулой $\chi_{{}_{\mathbb Q_\pad}}(a) = e^{2i\pi \{a\}_\pad}$, является аддитивным унитарным характером группы $(\mathbb Q_\pad, {+})$,
  причём каждый непрерывный унитарный  комплексный аддитивный характер  группы $(\mathbb Q_\pad, {+})$
  имеет вид $a\mapsto \chi_y(a) \equiv \chi_{{}_{\mathbb Q_\pad}}(y\cdot a)$ для некоторого ${y\in\mathbb Q_\pad}$, и~соответствие  $y\mapsto\chi_y$ является изоморфизмом  аддитивной топологической группы~$\mathbb Q_\pad$ и~(двойственной ей по Понтрягину) группы характеров на $\mathbb Q_\pad$.
  Та мера Хаара на локально компактной абелевой аддитивной группе $(\mathbb Q_\pad, {+})$, которая на каждом замкнутом шаре равна диаметру шара
  (совпадающему с~его радиусом),  называется далее канонической.


Под измеримостью отображений, определённых на измеримых пространствах и~принимающих значения в~нормированном поле $\pad$"=адических, вещественных или комплексных чисел % будем
понимаем измеримость относительно   сигма"=алгебры % $\sigma _\pad$
 всех борелевских подмножеств  в~соответствующем поле, обозначаемой $\sigma _\pad$, $\sigma _{\mathbb R}$ или $\sigma _{\mathbb C}$
 соответственно. %  $ \mathbb Q_\pad $ ~ . %

\smallskip

 \Section{Бесконечномерный оператор Лапласа  на $\pad$"=адическом пространстве
 и~соответствующее уравнение теплопроводности}
В этом пункте считается заданным вещественное число $T>0$, $Q$ "---  пространство  всех непрерывных справа функций  без разрывов второго рода,
 определённых на  вещественном отрезке  $[0; T]$  и~принимающих значения в~поле $\mathbb Q_\pad$.
 Далее, $P$ "--- пространство всех тех непрерывных слева функций  $[0; T]\to\mathbb Q_\pad$, каждая из которых  имеет % имеющих
 (свои) не более чем конечное число  точек разрыва и лишь конечное число значений.
 Зададим биекцию  между пространством $P$ и~ пространством всех тех $\mathbb Q_\pad$-значных борелевских мер на отрезке $[0; T]$,
 каждая из которых в~носителе имеет не более конечного числа точек,  формулой
 $p(s)=m_p[s; T]$  ($s\in[0; T]$), где $p\in P$ и $m_p$ "--- мера, отвечающая функции $p$.
Билинейное относительно поля $\mathbb Q_\pad $
 % \\
 отображение  
 $
 b :
 Q \times P  \to\mathbb Q_\pad
 $,
 задаваемое формулой 
 $$ 
 b (q,p) =-\int_{[0;T]} q(s)\;m_p(d\,s)
 =-\sum_{s\in\mathrm{supp\,} (m_p)}q(s)\cdot m_p(\{s\}),
$$
   приводит пространства
 $
 Q
 $
  и
 $
 P
 $
 в двойственность.



 Далее
 $\sigma _P$ "---
 сигма-алгебра с единицей $P$,
 порождённая всеми  функционалами вида
  $  b(q, {\cdot}) $
 ($q\in Q$);
 $M_P$ "--- пространство комплекснозначных  счётно"=аддитивных мер на  $(P,\sigma _P)$ и $F_Q$ "--- пространство их $b$"=преобразований Фурье, то есть  $b$-слабо секвенциально непрерывных  функций
 $
 \widetilde m : Q \to \mathbb C  $ вида
$$
 \widetilde m(q)=\int_P \chi_{{}_{\mathbb Q_\pad}}(b(q,p))\;m(d\,p),$$ где
 $m\in M_P$.
 Биекция $M_P\ni m\mapsto\widetilde m\in F_Q$ переносит норму, определенную в $M_P$
 полной вариацией мер, в банахову норму на $F_Q$.
 % Оператор
 Сходимость в пространстве $F_Q$ влечёт равномерную сходимость,
 так как $|\widetilde m(q)|\leqslant \|m\|_{M_P}=\|\widetilde m\|_{F_Q}$ \ ($m\in M_P$, $q\in Q$).
  Пространства $M_P$ и $F_Q$ не сепарабельны.
 Аналогично предыдущему определяем $\sigma _Q$
 как
   сигма-алгебру с единицей $P$,
   порожденную
 всеми функционалами вида
 $
  b({} \cdot {}, p)
 $
  ($p\in P$);
 $b$-преобразованием Фурье
 произвольной  счётно"=аддитивной меры
 $m:\sigma _Q\to\mathbb C$
 называем функцию $\widetilde m:P\to\mathbb C$,
 задаваемую формулой
 $\displaystyle \widetilde m(p)=\int_Q \chi_{{}_{\mathbb Q_{\pad}}}(b(q,p))m(dq)$.



 Если $p\in P$,
 то борелевским $p$-цилиндром (в пространстве $Q$)
  называется множество вида
 $
 C_{p,B}=\{q\in Q: b(q,p) % \varphi _p(q)
 \in B\}
 $,
 где $B\in\sigma _{\mathbb Q_\pad}$.
  Аналогично определяются борелевские $q$-цилиндры (в пространстве $P$).

\smallskip

\hypertarget{sham:th:1}{}

 \begin{theorem}[1] % {\bf Теорема 1.} \em
 Каждое из двух определенных выше $b$-преобразований Фурье
 является инъективным линейным оператором
 на своем пространстве мер.
 \end{theorem} % \rm

\smallskip

\begin{proof}
 Линейность очевидна.
  Для доказательства инъективности % достаточности
  можно провести рассуждения,
  аналогичные сделанным в~[\citen{shamarov:bib:SmShav}]\
  для числовых цилиндрических мер на вещественных
  пространствах,
  однако укажем
  более короткий путь,
   сформулировав его
  независимым от используемого поля ($\mathbb Q_\pad$ или ${\mathbb R}$)
  образом.
  Во-первых, в силу линейности достаточно проверить, что
  если преобразуемая мера $m$ ненулевая, то и её преобразование Фурье тоже.
  Пусть теперь $m\ne0$.
  Поскольку область определения меры
  $m$ как $\sigma $-алгебра порождена определенными выше борелевскими цилиндрами,
  мера принимает ненулевое значение на некотором цилиндре
  вида $\varphi ^{-1}( B)$,
  где $B$ "---
  борелевское   множество в поле, %$\mathbb Q_\pad$ 
   а
   $\varphi $ "---
  $b$-слабо непрерывный
  функционал на соответствующем пространстве.
  Таким образом, счётно"=аддитивная мера $m\circ \varphi ^{-1}$
  ненулевая на $\sigma _{\mathbb Q_\pad}$ или $\sigma _{\mathbb R}$ соответственно.
 Пространство счётно"=аддитивных мер,
 заданных на поле $\mathbb Q_\pad$ или ${\mathbb R}$,
  инъективно
 вложено в соответствующее пространство $S'$ обобщённых функций.
  Наконец, на  этом
 $S'$, как известно, преобразование Фурье --- автоморфизм,
 и, значит,
 $0\ne \widetilde{m\circ \varphi ^{-1}}\equiv\widetilde m\circ \varphi ^*$,
 где $\varphi ^*$ --- % алгебраически
 сопряженный к $\varphi $ оператор,
 откуда $\widetilde m\ne0$, что и требуется.
\end{proof}

\smallskip


 Фиксируем ещё вещественное число  $\alpha>0$  и~отвечающую ему неположительную вещественнозначную
 функцию
 $
 f
 :P\to(-\infty, 0]
 $,
 задав её формулой
 $ \displaystyle
 f
 (p)=-\int_0^T(|p(s)|_\pad)^\alpha ds,
  $
  где интеграл понимается в смысле Римана;
  функция  $f$
 измерима относительно $\sigma _P$.

\smallskip

 \begin{theorem}[2] % {\bf Теорема 1.} \em
 Для каждого $t\geqslant 0$ существует
 единственная
 вероятностная счётно"=аддитивная мера~%
 $w^t$ $($аналог меры Винера\/$)$
 на~$\sigma _Q,$
 обладающая преобразованием Фурье
% $$P^T\ni p\mapsto \int_Q\chi(b(q,p))\varphi ^s(dq)=\widetilde{\varphi ^t}(p)$$
 вида
 $$
 \displaystyle
 \widetilde{w^t}(p)=
 e^{t\cdot f(p)}=\exp \biggl({\displaystyle -t\cdot\int_0^T(|p(s)|_\pad)^\alpha ds }\biggr).
 $$
 \end{theorem} % \rm

\smallskip

\begin{proof}
  Для доказательства существования достаточно в качестве $w^t$
 взять распределение % построить
 однородного по времени и фазовому пространству $(\mathbb Q_\pad, {}+{})$
 (марковского) процесса $\xi^t$
 с~переходными плотностями вида
 $  p(t_1,t_2;x_1,x_2)  =     F_{(t_2-t_1)\cdot t}(x_2-x_1)  $,
где $0\leqslant  t_1<t_2\leqslant  T$, $x_1,$ $x_2\in\mathbb Q_\pad$,   и для каждого вещественного числа $s\geqslant 0$
и каждого $\pad$"=адического числа $x$  $
 F_s(x)
 =
 \sum_{k\in\mathbb Z}
               (e^{-s\pad^{k\cdot\alpha}}-e^{-s\pad^{(k+1)\cdot\alpha}})
              % \cdot
               \pad^{k}\Omega ( \pad^{-k}x ),
 $
 где, в свою очередь, $ \Omega (x)=1$ при $|x|_\pad\leqslant 1$
 и $\Omega(x)=0 $ в противном случае.
 При $t=1$ конструкция приведена в~[\citen{shamarov:bib:VVZ}],
 при других $t>0$ построения аналогичны,
 и, наконец, при $t=0$ вероятностная мера $m_1^0$
 сосредоточена в нуле пространства $P$.


  Единственность
 вытекает из доказанной в теореме~\hyperlink{sham:th:1}{1}  инъективности оператора
 преобразования Фурье.
\end{proof}

\smallskip 


 \begin{definition}
Пусть для произвольной  $ \sigma _Q$-измеримой  вещественной функции  $ f_0 : Q\to{\mathbb R} $, знакосочетание
 $(f_0 \cdot)$ означает  плотно определённый  (и,   вообще говоря, неограниченный)   оператор
  поточечного умножения 
  на
   (необязательно
  ограниченную)
  комплексную
 неположительную функцию $f_0 $  в банаховом пространстве  $M_P$,
 заданный
 на  (так называемой естественной) области определения  $D_{(f_0 \cdot)}=\{m\in M: (f_0 \cdot m)\in M_P  \}$  формулой
 $
 D_{(f_0 \cdot)}\ni m\mapsto {(f_0 \cdot m)}\in M_P.
 $
 Пусть  далее
 $\widehat {f_0 }$ --- оператор в $F_Q$  c~областью определения
 $
 D_{\widehat{ f_0  }}
 =
 \{ \widetilde m\in F_Q : m\in  D_{(f_0 \cdot)}  \}
 $,
 такой, что
 $ \widehat{ f_0  } (\widetilde m)= \widetilde{ f_0 \cdot m }$.
 \end{definition}


\smallskip

  Отметим, что
  подпространство
  $
  D_{(f_0 \cdot)}
    $
  плотно в банаховом $M_P$,
   так как $f_0$ всюду конечна и  $\sigma _P$-измерима,
  и что в случае ограниченного сверху множества значений
  функции $f_0$,
  $(f_0 \cdot)$ "--- генератор сильно непрерывной полугруппы
  $\{(e^{t\cdot f_0}\cdot): t\geqslant 0\}$
  операторов в банаховом пространстве $M_P$.

\smallskip

\hypertarget{sham:lemma:1}{}
 \begin{lemma}[1]
   В пространстве $M_P$
  корректно определён
  билинейный оператор свёртки мер,
  переводящий упорядоченную пару мер
  в образ
  их прямого произведения
  относительно сложения.
  Преобразование Фурье переводит такую свёртку мер
  в произведение их преобразований Фурье.
 \end{lemma}

\smallskip
{\it   Д\,о\,к\,а\,з\,а\,т\,е\,л\,ь\,с\,т\,в\,о} использует измеримость операции сложения
  в $P$ относительно алгебры борелевских $b$-цилиндров,
  определяемых как прообразы борелевских множеств % из $\mathbb Q_\pad^n$
  относительно $b$-слабо непрерывных
  конечномерных линейных отображений на $P$.

\smallskip
 \begin{theorem}[3]
  Оператор $  \widehat{f } $
 является генератором
 сильно непрерывной
 свёрточной операторной  полугруппы $\{(w^t*): t\geqslant 0\}$
 в пространстве $F_Q.$
 \end{theorem}
 
 \smallskip


 Эта теорема есть следствие
 корректности определения свёртки,
 двух пре\-д\-ыдущих теорем, леммы \hyperlink{sham:lemma:1}{1} и того факта, что
 $(f\cdot)$ --- генератор
 однопараметрической
 сильно непрерывой полугруппы
 $\{e^{t\cdot(f\cdot)}\}_{t\geqslant 0}$
 сжатий
 в банаховом % пространстве
 $M_P$.
 

 \begin{definition}
 Оператор
 $
 \widehat{f} % =\widehat{f_\alpha }
 $,
 плотно определённый  в банаховом пространстве $F_Q$,
  далее называется  бесконечномерным оператором Лапласа  в~степени $\alpha/2$,
 причём в случае $\alpha=2$ " --- просто  бесконечномерным оператором Лапласа.
 \end{definition}

\smallskip



 Поставим задачу Коши
  с соответствующим
  бесконечномерному оператору Лапласа
 уравнением теплопроводности
 (с включенным распределённым источником).
 Фиксируем меру $m\in M_P$.
 
 \smallskip

 \begin{definition}
 Пусть   $\psi_0\in F_Q$ "--- начальное условие задачи, необходимо  найти  непрерывную  функцию
 $\Psi:[0,\infty)\to F_Q$,
  дифференцируемую  при каждом $t>0$  и такую,  что
\begin{equation}
\label{sham:eq1}
 \Psi'(t)=\widehat {f }\; \Psi_t+ (\widetilde{m }\cdot) \Psi_t
 \quad (t>0), \quad \Psi(0)=\psi_0,
\end{equation}
где $\Psi_t=\Psi(t)$ ($t\in[0; +\infty)$).
 Эту задачу будем называть задачей \eqref{sham:eq1}.

 Далее при $t\geqslant 0$ и $q\in Q$  запись $\psi(t,q)$ означает
  $  \Psi_t(q) $.
 \end{definition}

 \smallskip
 

 \Section{Формулы Фейнмана и Фейнмана"--~Каца для задачи (\ref{sham:eq1})}
Сумма генератора   $\hat f$  и ограниченного оператора $ (\widetilde{m
}\cdot) $  сама является генератором сильно непрерывной полугруппы,
 элементы которой обозначаются $ e^{t\cdot(\hat f+ (\widetilde{m }\cdot) )}$.
 При этом справедлива  следующая теорема (формула Троттера).

\smallskip
 \begin{theorem}[4]
  Для всякого $\psi_0\in F$ 
 $$
 \Psi(t)= e^{t\cdot\left(\widehat{ f }+ (\widetilde{m }\cdot) \right)}\psi_0
 =\lim_{n\to\infty}(e^{(t/n)\cdot \widehat{ f } }
 e^{(t/n)\cdot(\widetilde{m }\cdot) })^n\psi_0.
 $$
 \end{theorem}
\smallskip

\begin{lemma}[1]
 Для всякого $\psi_0\in F$ и $ q\in Q $
 справедливы равенства
\begin{multline*}
 \psi(t,q)=
  \lim_{n\to\infty}
  \bigl((e^{(t/n)\cdot \widehat{ f_0} }
  e^{(t/n)\cdot(\widetilde{m }\cdot) })^n
  \psi_0\bigr)(q)
 =
\\
 =\lim_{n\to\infty}
   \left(
   \bigl((\omega ^{t/n }\;*)\circ
 (e^{(t/n)\widetilde{m }}\;\cdot) \bigr)^n
 \ \psi_0\right)
 (q).
\end{multline*}
  \end{lemma}


\begin{newthm}{Следствие}
  Положим
 $I_n(q)=
    \left(
   ((\omega ^{t/n }\;*)\circ
 (e^{(t/n)\widetilde{m }}\;\cdot) )^n
 \ \psi_0\right)
 (q)
=$
%
%
 $$
  =
 \int \omega ^{t/n }(q-dq_{n-1}) e^{(t/n)\widetilde{m }(dq_{n-1})}
 \int \omega ^{t/n }(q_{n-1}-dq_{n-2}) e^{(t/n)\widetilde{m }(dq_{n-2})}
 \int
 \cdots
 $$
 $$
 \dots
 \int \omega ^{t/n }(q_{2}-dq_{1}) e^{(t/n)\widetilde{m }(dq_{1})}
 \int \omega ^{t/n }(q_{1}-dq_{0}) e^{(t/n)\widetilde{m }(dq_{0})}
  \psi_0(q_0)
  =
  $$
 $($замена $q_j=q-\widetilde q_j)$
 $$
  =
 \int \omega ^{t/n }(d\widetilde q_{n-1}) e^{(t/n)\widetilde{m }(q-\widetilde q_{n-1})}
 \int \omega ^{t/n }(-\widetilde q_{n-1}+d\widetilde q_{n-2}) e^{(t/n)\widetilde{m }(q-\widetilde q_{n-2})}
 \int
 \cdots
 $$
 $$
 \dots
 \int \omega ^{t/n }(-\widetilde q_{2}+d\widetilde q_{1}) e^{(t/n)\widetilde{m }(q-\widetilde q_{1})}
 \int \omega ^{t/n }(-\widetilde q_{1}+d\widetilde q_{0}) e^{(t/n)\widetilde{m }(q-\widetilde q_{0})}
  \psi_0(q-\widetilde q_0).
  $$
  Тогда
  $
  \psi(t,x)
  =
  \lim_{n\to\infty}
  I_n(q).
  $
 \end{newthm}

\hypertarget{def:shamarov:i}{}

 \begin{definition}[1]Пусть $t>0$,  пространство $\Gamma^t_Q$ состоит из  всех функций \mbox{$\gamma:[0; t]\to Q $,}  не имеющих разрывов второго рода,
 и~пусть $G$ и $g$ "--- ограниченные непрерывные функции  $Q\to\mathbb C$.
 Интегралом $$\int_{\Gamma^t_Q} \exp\biggl({\int_0^t g(\gamma(s))ds}\biggr)\cdot G(\gamma(0)) M^t_\omega (d\gamma)$$ % Фейнмана
 по центральной псевдомере $M^t_\omega $,
 определяемой
 переходными мерами 
 $$\bigl\{\omega ^{s}{(A{-}q)}:s\geqslant 0,\quad q\in Q,\quad A\in \sigma _Q\bigr\},$$
 от функции 
 $$
 \Gamma\ni\gamma \mapsto \exp \biggl({\int_0^t g(\gamma(s))ds} \biggr) \cdot G(\gamma(0))
 $$
 назовём предел при $n\to\infty$  интегралов вида
  $$
 \int \omega ^{t/n }(d q_{n-1})e^{(t/n)g( q_{n-1})} 
 \int \omega ^{t/n }(-  q_{n-1}+d q_{n-2})e^{(t/n)g( q_{n-2})}
 \int
 \cdots
 $$
 $$
 \dots
 \int \omega ^{t/n }(-  q_{2}+d  q_{1})e^{(t/n)g( q_{1})}
 \int \omega ^{t/n }(- q_{1}+d q_{0})e^{(t/n)g(  q_{0})}
  G ( q_0).
  $$
   \end{definition}

\smallskip
   \begin{lemma}
  Определение~\hyperlink{def:shamarov:i}{\rm 1} корректно.
   \end{lemma}

\smallskip

   \begin{newthm}{Теорема 5 (Формула Фейнмана"--~Каца)}
  Для всякого $\psi_0\in F$ и $ q\in Q $
 справедливо равенство
 $$
 \psi(t,x)
 =
 \int_{\Gamma ^t_Q} \exp \biggl({\int_0^t \widetilde {m_0}(q-\gamma(s))ds}\biggr) \cdot \psi_0(q-\gamma(0))
 M^t_\omega (d\gamma).
 $$
   \end{newthm}


\smallskip



 {\bf Заключение.}   Одним из ключевых технических средств в конструкции
 разрешающей полугруппы для эволюционного уравнения
 является теорема о формуле Троттера, которая в~настоящей работе
 применена для получения  формул, называемых формулами Фейнмана и Фейнмана"--~Каца в
 конфигурационном пространстве. Возможно,  псевдомеры $M^t_\omega $ счётно"=аддитивны.



 \thanks{Работа выполнена при поддержке РФФИ (проект \No~10--01--00724--a).}

%\newpage
%


\begin{thebibliography}{9}


 \RBibitem{shamarov:bib:VVZ}
\by Владимиров~В.\,С, Волович~И.\,В., Зеленов~Е.\,И.
\book $p$-Адический анализ и математическая физика
\yr 1994
\publ Физматлит
\publaddr М.
\totalpages 352
\transl
англ. пер.:~
\by Vladimirov~V.\,S., Volovich~I.\,V., Zelenov~E.\,I.
\book $p$-Adic analysis and mathematical physics
\serial Series on Soviet and East European Mathematics
\yr 1994
\publ World Scientific Publishing Co., Inc.
\publaddr River Edge, NJ
\vol 1
\totalpages 319

\RBibitem{shamarov:bib:SSK}
\by   Смолянов~О.\,Г.,  Шамаров~Н.\,Н.,  Кпекпасси~М.
\paper  Формулы Фейнмана"--~Каца и Фейнмана для бесконечномерных уравнений с~оператором Владимирова
\jour  Докл. РАН
\yr 2011
\vol 438
\toappear
\misctransl 
\jour Dokl. RAN
\yr 2011
\paper Feynman and Feynman--Kac formulas for infinite-dimensional equations with Vladimirov operator
\by Smolyanov~O.\,G.,  Shamarov~N.\,N., Kpekpassi~M.
\yr 2011
\vol 438
\toappear


 \Bibitem{shamarov:bib:BeloshQN}
 \by Beloshapka~O.\,V.
 \paper Feynman formulas for an infinite dimensional $p$-adic heat type equation
 \jour  IDAQP
 \yr 2011
 \vol 14
 \issue 1
\pages 137--148
%DOI: 10.1142/S021902571100433X


\RBibitem{shamarov:bib:Stekl}
\by Смолянов~О.Г., Шамаров~Н.Н.
\paper Формулы Фейнмана и~интегралы по траекториям для эволюционных уравнений с~оператором Владимирова
\inbook Избранные вопросы математической физики и~$p$-адического анализа
\bookinfo Сборник статей
\serial Тр. МИАН
\yr 2009
\vol 265
\pages 229--240
\publ МАИК
\publaddr М.
\transl англ. пер.:~
\paper Feynman formulas and path integrals for evolution equations with the Vladimirov operator
\by  Smolyanov~O.\,G.,  Shamarov~N.\,N.
\jour Proc. Steklov Inst. Math.
\yr 2009
\vol 265
\pages 217--228






 
\RBibitem{shamarov:bib:S82}
\by   Смолянов~О.\,Г.
\paper  Бесконечномерные псевдодифференциальные операторы и квантование по Шредингеру
\jour  Докл. Акад. наук СССР
\yr 1982
\vol 263
\issue 3
\pages 558--562
\transl англ. пер.:~
\by Smolyanov~O.\,G.
\paper Infinite-dimensional pseudodifferential operators and Schr\"odinger quantization
\jour Sov. Math. Dokl.
\vol 25
\issue 3
\pages 404--408 
\yr 1982


\RBibitem{shamarov:bib:AS}
\by Аккарди~Л.,  Смолянов~О.\,Г.
\paper  Формулы Фейнмана для эволюционных уравнений с~лапласианом Леви на бесконечномерных многообразиях
\jour Докл. РАН
\yr 2006
\vol 407
\issue 5
\pages 583--588
 \transl англ. пер.:~
\jour  Doklady Mathematics
\vol 73
\issue 2
\pages 252--257
%DOI: 10.1134/S106456240602027X
\paper Feynman formulas for evolution equations with Levy Laplacians on infinite-dimensional manifolds
\by Accardi~L., Smolyanov~O.\,G.
\yr 2006

\RBibitem{shamarov:bib:SmShav}
\by Смолянов~О.Г., Шавгулидзе~Е.Т.
\book Континуальные интегралы
\publaddr  М.
\publ МГУ
\yr 1990
\totalpages 150
\misctransl
\by  Smolyanov~O.\,G.  Shavgulidze~E.\,T.
\book Continual Integrals
\publaddr  Moscow
\publ MGU
\yr 1990
\totalpages 150


% \Bibitem{shamarov:bib:FN}
%\by  Feynman~R.\,P.
%\paper The development of the space-time view of quantum electrodynamics
%\jour Science 
%\vol 153 
%\yr 1966
%\issue 3737 
%\pages 699--708
%%DOI: 10.1126/science.153.3737.699
%\rtransl русск. пер.:~ 
% \by Фейнман Р.
% \paper Развитие пространственно"=временн\'ой трактовки квантовой электродинамики (нобелевские лекции по физике 1965~г.)
% \jour  УФН
% \yr 1967
% \vol 91
% \issue 1
% \pages 29--48
% 
% 
%\Bibitem{shamarov:bib:F}
%\by Feynman~R.\,P.
%\paper Space-time approach to non-relativistic quantum mechanics
%\jour Rev. Mod. Phys.
%\vol 20
%\issue 2
%\pages 367–387 
%\yr 1948
%
%
%
% \RBibitem{shamarov:bib:Dynkin}
% \by Дынкин~Е.\,Б.
% \book Основания теории марковских процессов
% \publaddr  М.
% \publ Физматгиз
% \yr 1959
% \totalpages 227
% \misctransl 
% \by Dynkin~E\,B.
% \book Foundations of the Theory of Markov Processes
% \publaddr Moscow
% \publ Fizmatgiz
% \yr 1959
% \totalpages 227
%
%
%
%
% \RBibitem{shamarov:bib:Bikul}
% \by Аветисов В.А., Бикулов А.Х., Осипов В.Ал.
%\paper $p$-Адические модели ультраметрической диффузии в~конформационной динамике макромолекул
%\inbook Избранные вопросы $p$-адической математической физики и анализа
%\bookinfo Сборник статей. К 80-летию со дня рождения академика Василия Сергеевича Владимирова
%\serial Тр. МИАН
%\yr 2004
%\vol 245
%\pages 55--64
%\publ Наука
%\publaddr М.
%\transl англ. пер.:~
%\by Avetisov~V.\,A., Bikulov~A.\,Kh., Osipov~V.\,A.
%\paper $p$-Adic models for ultrametric diffusion in conformational dynamics of macromolecules
%\jour Proc. Steklov Inst. Math.
%\yr 2004
%\vol 245
%\pages 48--57
%
%
%\Bibitem{shamarov:bib:Shaef}
%\by Schaefer~H.\,H.
%\book Topological vector spaces
%\serial  Macmillan series in advanced mathematics and theoretical physics
%\totalpages 294
%\publ Macmillan 
%\publaddr New York
%\yr 1966  
%\rtransl русск. пер.:~
%\by Шефер~Х.
%\book Топологические векторные пространства
%\publaddr  М.
%\publ Мир
%\yr 1975
%\totalpages 359 
%
%
%\RBibitem{shamarov:bib:Chernov}
%\by Chernoff~P.\,R.
%\paper A note on products formulas for operator semigroups
%\jour  J.~Functional Anal.
%\yr 1968
%\vol 2
%\issue  2
%\pages 238--242
%
%
%
%
%
%\Bibitem{shamarov:bib:SmTkTr}
% \by Smolyanov~O.\,G., Tokarev~A.\,G., Truman~A. 
% \paper Hamiltonian Feynman path integrals via the Chernoff formula 
% \jour  J. Math. Phys.
% \yr 2002
% \vol 43
% \issue 10
% \pages 5161--5171






 \end{thebibliography}
 
 

\hskip6.5cm \thedate \par
\hskip6.5cm\therevdate\par

\newpage 

\chead[\fancyplain{}{\scriptsize \sl S\,h\,a\,m\,a\,r\,o\,v~N.\,N.}]{\fancyplain{}{\scriptsize \sl  Functional Laplace operator on a  $\pad$"=adic space \dots}}


\makeengtitlem

\renewcommand{\baselinestretch}{.9}

\newpage

%\end{document}
%
